\documentclass[12pt]{article}
\title{Trade Theory: Begin Here}
\date{Semester 02, 2026}
\author{Nikhil Damodaran}
\usepackage{/Users/ndamodaran/Library/CloudStorage/Dropbox/Nikhil/bibliography/article-nikhil}


\begin{document}

\maketitle


This is a reading list mostly on classical trade theories, trade policy and some recent developments in the field within these sub fields. Generally, these topics are less interesting to the journals so there needs to be a newer empirical or theoretical hook for work on these models to continue. Or the more recent papers deal with contemporary questions and give answers to the discussion of our times.

\section{Comparative Advantage}
This is probably the most important concept in trade and the most basic. To say that most of trade is a result of technological differences because scarcity rules is to say something basic. However, it is a good starting point in analysis.

\subsection{Pedagogy}
\cite{Findlay1987D}

\cite{Bernhofen2018JEP}

\subsection{Theory}
-- \cite{Dornbusch1977AER}

The definitive many good version of Ricardian model which has become the textbook version of Ricardian model at higher levels.

-- \cite{Costinot2015QJE}

This is a very important paper. The idea is that while we have a reasonable understanding of comparative advantage, we do not have an understanding of what this model says about trade policy and intervention. This is an attempt at answering different set of questions using the model which posits comparative advantage. * This is an interesting way of posing a question too. If we cannot contribute by asking old questions, ask new questions from the old model.

-- \cite{Gaubert2021JPE}

This is the most recent work using comparative advantage. This also could be the intersection of international trade (a conventional field) and industrial policy which specifically focuses on how to curb the market power of some large firms. They break down comparative advantage to an average firm behavior v. behavior generated by some large firms.

-- \cite{Grossman1990AER}: Well cited

-- \cite{Bernard2007RES}: Very well cited. The theory combines a two-country model with specialization based trade with heterogeneous firms.

-- \cite{Clarida1992AERPP}

-- \cite{Costinot2009E}

-- \cite{Deardorff1980JPE}

-- \cite{Deardorff2005RIE}

-- \cite{Costinot2009JIE}

-- \cite{Neary2007RES}

-- \cite{Keesing1966AER}

-- \cite{Redding1999OEP}

-- \cite{Matsuyama1992JET}

\subsection{Empirics}

-- \cite{Bernhofen2004JPE}

-- \cite{Costinot2012AERPP}

-- \cite{Levchenko2016JME}

A sectoral take on specialization and comparative advantage. Cross country regressions but at the sectoral level. Thinking about doing this work, if it could be done, seems silly at this stage. Especially because at the time this work was published, cross country convergence regressions were discarded and empirical work would have been looked down upon.

-- \cite{Golub2000RIE}

-- \cite{Chor2010JIE}

-- \cite{Hummels1993AERPP}

-- \cite{Hoen2006ARS}


\subsection{Comparative Advantage Applied Elsewhere}

-- \cite{Boehm2022JPE}

Interesting piece of evidence on Indian firms production structure. The idea is that if firms have readily available inputs, then they tend to specialize and produce in the sector which has fewer input frictions.

-- \cite{Davis2020JIE}

The idea is that of underlying labor specialization which causes some cities to specialize in some products. The intersection of labor and trade is very nice.

-- \cite{Hunt1995JM}

This piece is even more cited than the benchmark comparative advantage work by DFS. It is applying comparative advantage as a concept to marketing and explaining the monopoly power of firms. 5K citations and counting.

-- \cite{Suri2011E}

This is the concept adapted to technology adoption by farmers. The data applied to an econometric model is what makes this work worthy of econometrica. It is a sharp neat application.

-- \cite{Nunn2014HIE}

This paper examines how domestic institutions can be a source of comparative advantage. So rather than pinning comparative advantage to technology or labor productivity, this talks about institutions such as what Adam Smith mentions when he talks about institutions supporting free market to function and trade to occur.

-- \cite{Liu2020JDE}

-- \cite{Do2016JDE}

-- \cite{Ju2011JIE}

\section{Exchange Rates}

\subsection{Purchasing Power Parity}
The idea here is that purchasing power parity is a long run phenomenon with a convergence to this anchor in the very long run. Rogoff gives the 3-5 year half life result. Since the convergence is so slow, people are less interested in this theoretical long term benchmark but rather interested in short run exchange rate movements.

A leading reason for this lack of interest is: rather than viewing exchange rate parity as a result of equalization of goods prices, it started being viewed as a result of financial interaction, at least in the short run. This implies that things such as monetary policy differences across countries, capital flows, international exchange rate futures trade became more important to the determination of exchange rate.

A second reason for the disconnect is that firms price discriminate across countries which implies that exchange rate pass-through is incomplete and local currency pricing matters. Hence if by construction firms are treating different international markets differently, then the difference in exchange rate will reflect the difference in import/domestic consumption composition.

A third reason for the disconnect is that a large part of domestic prices is explained by non-tradables such as rents, local services such as haircuts, restaurants etc. Hence, when the CPI has more non-tradables, there is no reason to assume PPP equalization. This is connected to the idea extended by Balassa-Samuelson when they argued that countries will differ in price levels because of differences in productivity across tradeables which permeate into non-tradables.

\cite{Rogoff1996JEL}

\cite{Balassa1964JPE}

\cite{Pinkovskiy2020AEJM}

\cite{Taylor2002JEP}

\cite{Dornbusch1987PDE}

\cite{Vo2023JES}

\cite{Frenkel1978JIE}

\section{Tariffs}

\subsection{Empirical Results}
\cite{Suwanprasert2026IE}


















\clearpage
\begingroup
\footnotesize
\setlength{\bibsep}{0.2em}
\bibliographystyle{apalike}
\bibliography{/Users/ndamodaran/Library/CloudStorage/Dropbox/Nikhil/bibliography/References}
\endgroup
\end{document}